\section{LQG：看不全状态时怎样继续反馈}
\label{sec:16-Control-and-Reinforcement-Learning/01-Optimal-Control/05}

一间小机房要控送风温度，可探头只有一个，装在回风口。机柜的发热随业务负载起伏，墙体和地板还在慢慢蓄放热量，探头本身有正负零点几度的抖动。上一节的 LQR 张口就要完整状态：每排机柜的温度、蓄热体的当前热量、扰动的大小。你手上是一个带噪声的标量读数。把读数直接当状态塞给 \(u=-Kx\) 会怎样？噪声被增益放大打进风机，蓄热那部分状态压根没进控制律。反馈还想继续做下去，中间得补一块东西。

\subsection{读数不是状态，中间隔着一层噪声}

把场景写成标准形式。状态 \(x_t\) 按线性动力学演化并受随机扰动，观测是状态的线性函数再加噪声：

\[
x_{t+1} = A x_t + B u_t + w_t, \qquad y_t = C x_t + v_t,
\]

其中 \(w_t \sim \mathcal{N}(0, W)\) 是过程噪声，\(v_t \sim \mathcal{N}(0, V)\) 是测量噪声，两者独立且各时刻独立。控制器每一步能看到的只有历史观测 \(y_0,\dots,y_t\) 和自己下过的历史控制。

这样一来控制律的定义域变了。LQR 的策略是状态的函数，现在能用的只有观测历史的函数。策略空间大得多也乱得多——一个可以任意读取全部历史的非线性映射，怎么找最优的那个？

线性二次高斯（LQG）问题就是在这个设定下，最小化

\[
\E \sum_{t} \bigl( x_t^{\trans} Q x_t + u_t^{\trans} R u_t \bigr),
\qquad Q \succeq 0,\; R \succ 0 .
\]

代价里出现的仍然是真实状态 \(x_t\)，尽管你看不见它。

\subsection{Kalman 滤波器把读数递推成条件均值}

线性加高斯有一个便利的后果：在给定观测历史的条件下，\(x_t\) 的后验分布仍是高斯，因此被均值与协方差两个量完整刻画，不必存整段历史。递推分预测与更新两步：预测把上一步的后验推过动力学并加上 \(W\)，不确定性只增不减；更新用新息（innovation），也就是读数与预测读数之差，把它拉回来：

\[
r_t = y_t - C\widehat x_{t\given t-1},
\qquad
K_t = P_{t\given t-1} C^{\trans}
\bigl( C P_{t\given t-1} C^{\trans} + V \bigr)^{-1},
\qquad
\widehat x_{t\given t} = \widehat x_{t\given t-1} + K_t r_t .
\]

增益为什么恰好是这个形状、它怎样随预测不确定性与探头噪声的比值摆动、以及整条闭式递推在什么地方断裂，上一节 \hyperref[sec:16-Control-and-Reinforcement-Learning/01-Optimal-Control/04]{``Kalman 滤波：用预测与更新的循环从噪声中估计状态''} 已经推完；从联合分布的条件独立结构进入同一件事，则见 \hyperref[sec:13-Machine-Learning/07-Graphical-and-Sequential-Models/03]{``Kalman filter 是动态系统上的 Gaussian 条件化''}。本节只取这套递推的输出 \(\widehat x_{t\given t}\)，把它接到控制上。

接之前先记下一条性质：误差协方差的递推里没有出现 \(y_t\)，它的演化只取决于 \(A, C, W, V\)，可以离线算完存表。这条性质等下要派用场。

\subsection{分离定理让两个 Riccati 方程各解各的}

现在把估计接到控制上。直觉上的做法是拿估计值冒充真状态，套用上一节的 LQR 增益。这种做法叫确定性等价（certainty equivalence），听上去像是偷懒的近似。在 LQG 这组假设下，它恰好是最优的。

\begin{theorem}[分离定理]
设动力学线性、代价二次、噪声高斯且与控制独立，模型参数已知，系统可控可观。则在所有以观测历史为输入的策略中，最优控制律为
\[
u_t = -L_t\, \widehat x_{t\given t},
\]
其中 \(L_t\) 是完整状态可观测时 LQR 的反馈增益，\(\widehat x_{t\given t}\) 是 Kalman 滤波器给出的条件均值。
\end{theorem}

证明的骨架是一次配方。把代价里的 \(x_t\) 拆成 \(\widehat x_{t\given t} + e_t\)，其中 \(e_t = x_t - \widehat x_{t\given t}\) 是估计误差。高斯条件下 \(e_t\) 与观测历史正交，因而与任何一个合法控制 \(u_t\) 不相关，交叉项在期望下消失。总代价于是裂成两块：一块只含 \(\widehat x\) 与 \(u\)，形状与完整状态的 LQR 一模一样；另一块只含 \(e_t\) 的协方差，而这个协方差按上一小节的观察压根不受控制影响。第二块是一个加性常数，最小化时可以整个丢掉，第一块的最优解就是 LQR 增益作用在 \(\widehat x\) 上。

条件对账要记清楚。交叉项消失用的是高斯下``条件均值即正交投影''，非高斯时只剩线性估计的正交性，配方不再干净；误差协方差与控制无关用的是观测方程 \(y_t = Cx_t + v_t\) 里控制不影响噪声强度，也不影响 \(C\)；\(R\succ 0\) 与可控可观保证 LQR 那一块的 Riccati 方程有稳定解。

结论的工程含义是两个 Riccati 方程可以彻底分家。控制侧那个从终端往回积，只吃 \(A, B, Q, R\)；估计侧那个往前推，只吃 \(A, C, W, V\)。调噪声模型不必重算增益，调代价权重不必重跑滤波器。

\begin{center}
\begin{tikzpicture}[>=stealth,
  block/.style={draw,rounded corners=2pt,minimum width=2.3cm,minimum height=0.95cm,fill=blue!6},
  every node/.style={font=\small}]

  \draw[dashed,gray!70] (-1.45,-0.85) rectangle (1.45,0.85);
  \node[gray!55!black,font=\footnotesize] at (0,1.15) {确定性最优反馈};
  \draw[dashed,gray!70] (1.65,-2.65) rectangle (4.55,-1.0);
  \node[gray!55!black,font=\footnotesize] at (3.1,-2.95) {最优状态估计};

  \node[block] (gain) at (0,0) {增益 \(-L_t\)};
  \node[block] (plant) at (4.1,0) {被控对象};
  \node[block] (kf) at (3.1,-1.85) {Kalman 滤波器};

  \draw[->,thick] (gain) -- node[above] {\(u_t\)} (plant);
  \draw[->,thick] (plant.east) -- (6.9,0);
  \node[anchor=south] at (6.5,0.05) {\(y_t\)};
  \draw[->,thick] (6.9,0) -- (6.9,-1.85) -- (kf.east);
  \draw[->,thick] (kf.west) -- (-2.1,-1.85) -- (-2.1,0) -- (gain.west);
  \node[anchor=south] at (-0.6,-1.8) {\(\widehat x_{t\given t}\)};

  \draw[->,gray!80] (4.1,1.3) -- (plant.north);
  \node[gray!55!black,anchor=west] at (4.2,1.15) {\(w_t\)};
  \draw[->,gray!80] (5.9,1.3) -- (5.9,0.03);
  \node[gray!55!black,anchor=west] at (6.0,1.15) {\(v_t\)};

  \draw[->,gray!60,dashed] (2.05,-0.15) -- (2.75,-1.35)
    node[midway,anchor=west,gray!55!black,font=\footnotesize] {\(u_t\)};
\end{tikzpicture}

\emph{两个虚线框各自独立设计：左框只看 \(A,B,Q,R\)，右框只看 \(A,C,W,V\)；它们靠 \(\widehat x\) 和 \(u\) 串成一个闭环。}
\end{center}

\subsection{分开设计不等于互不影响}

上图容易读出一个过头的结论：既然两个框各管各的，估计做得好不好就与控制表现无关。并非如此。

设计过程分家，运行时两者仍在同一个回路里。\(W\) 与 \(V\) 决定了误差协方差的稳态水平，误差直接进入实际代价的那个加性常数——分离定理说它不影响最优增益的选择，没说它不影响你最终付多少钱。探头换成更差的，最优控制律一个字不改，账单照涨。图里那条从 \(u_t\) 指向滤波器的虚线也提醒着另一层耦合：控制值本身是预测步的输入，估计器必须知道你下了什么指令，否则 \(B u_{t-1}\) 这一项对不上，新息会系统性偏移。

分离定理省掉的是搜索。它证明了不必在``所有观测历史的非线性函数''这个庞大空间里找最优策略，那个最优策略恰好落在``先估计再线性反馈''这个狭窄的子集里。省掉的是搜索范围，不是耦合本身。

\subsection{离开线性高斯，确定性等价就开始漏水}

分离定理的四个前提每一个都能被现实拿走，拿走的方式也各不相同。

动力学非线性时，后验分布不再是高斯，条件均值一般不能用有限个数递推。扩展 Kalman 滤波器在当前估计处线性化，粒子滤波用带权样本近似后验，两者都提供估计，但没有哪一个还配着``估计加 LQR 增益即最优''的保证。它们是可用的工程方案，不是分离定理的推论。

噪声非高斯时，条件均值仍是均方误差意义下的最优估计，但计算它一般要跟踪整个后验；退而求其次用线性最小方差估计，就得接受由此产生的次优性。

约束一旦激活，配方那一步就塌了。带饱和的最优控制律不是线性的，把估计值代进一个非线性律，误差项与控制项的交叉不再消失。

最深的一处漏水叫对偶效应（dual effect）。上面用到``误差协方差不受控制影响''，这条性质依赖观测方程里控制不改变可观测性。系统参数未知、或者观测质量本身取决于状态时，动作会改变你未来能看清多少。此时最优策略会主动做出一些当下看似次优的动作，只为把系统推到信息量更大的区域去——探索与利用的权衡在这里第一次出现，而确定性等价对它完全无感。这条线索会在 \hyperref[ch:136]{``强化学习：模型未知时怎样逼近最优决策''} 里成为主角。

回到机房。如果墙体蓄热被建模得太弱，滤波器算出的协方差可能看起来很小，房间却持续超调——估计器对自己的错误很有信心。这类故障有一个统一的检查口：新息序列。模型与噪声设定正确时，\(r_t\) 应当是零均值、互不相关、协方差等于 \(CPC^{\trans}+V\) 的白噪声。它若有偏、有自相关，或者实际方差与理论值对不上，说明该修的是模型，不是增益。LQG 给出的是一条能从头推到尾的基准线，它把估计、反馈与不确定性怎样严丝合缝地接在一起说清楚了；超出线性高斯的部分，得靠别的工具接着走。
